\section{Modelado e identificación del sistema}

\subsection{Modelo físico}

El modelo se planteó desde primeros principios, con las ecuaciones de Lagrange para la bola
sobre el plato. Las hipótesis fueron: la bola rueda sin deslizar y no gira alrededor de la normal
del plato; es rígida, con momento de inercia $J = k\,m\,r_b^2$; el plato y el mecanismo son
rígidos y sin holgura; no hay fricción de rodadura; la entrada del modelo es la inclinación del
plato, con el actuador modelado aparte; y la cámara está fija sobre el eje de la junta. Se define
el coeficiente de rodadura $c = 1/(1+k)$, que vale $5/7$ para una esfera maciza. La masa de la
bola se cancela.

Para un eje, con el otro horizontal, la ecuación queda

\begin{equation}
\ddot x = c\left[-g\sin\theta + x\,\dot\theta^{2} + \hef\,\ddot\theta\right],
\qquad \hef = h + k\,r_b,
\label{eq:nolineal}
\end{equation}

donde $h$ es la altura del centro de la bola sobre el centro de la junta. El modelo completo de
dos ejes agrega términos de Coriolis y de acople que solo aparecen cuando los dos ejes se mueven
a la vez. Las no linealidades son el seno, el término centrífugo $x\,\dot\theta^{2}$, la
proyección que hace la cámara cuando el plato está inclinado, el mecanismo de cuatro barras
entre servo y plato, y la saturación y la velocidad máxima del servo.

La cámara no mide la posición sobre el plato. Como la homografía se calibró con el plato
horizontal, con el plato inclinado mide la proyección del centro de la bola, que linealizada da
$\Xm = x - h\,\theta$. Linealizando \eqref{eq:nolineal} alrededor de $x = \dot x = 0$,
$\theta = \dot\theta = 0$, se obtienen dos funciones de transferencia por eje,

\begin{align}
\frac{X(s)}{\Theta(s)} &= \frac{c\,\hef\,s^{2} - c\,g}{s^{2}},
\label{eq:lineal}\\
\frac{\Xm(s)}{\Theta(s)} &= \frac{-(1-c)\,\hsup\,s^{2} - c\,g}{s^{2}},
\label{eq:linealm}
\end{align}

donde $\hsup$ es la altura de la cara del plato sobre la junta. El jacobiano del modelo de dos
ejes en el punto de operación no tiene términos cruzados (diferencia relativa con el analítico de
\num{1.7e-13}), así que a primer orden los ejes quedan desacoplados y se tratan como dos lazos
SISO iguales. La Tabla~\ref{tab:fisico} resume los valores con la geometría del diseño.

\begin{table}[!t]
\caption{Modelo lineal por eje con los parámetros del diseño ($c = 5/7$,
$\hsup = \SI{49}{\milli\metre}$, $h = \SI{69}{\milli\metre}$, $r_b = \SI{20}{\milli\metre}$).
Fuente: \texttt{etapaA/resultados/A2\_consola.txt}.}
\label{tab:fisico}
\centering
\footnotesize
\begin{tabular}{@{}lr@{}}
\toprule
\textbf{Magnitud} & \textbf{Valor} \\
\midrule
$c\,g$ & \SI{7007}{\milli\metre\per\second\squared\per\radian} \\
$c\,\hef$ & \SI{55.0}{\milli\metre\per\radian} \\
$-(1-c)\,\hsup$ & \SI{-14.0}{\milli\metre\per\radian} \\
Polos & $s = 0$ (doble) \\
Ceros de $X/\Theta$ & $s = \pm\SI{11.29}{\radian\per\second}$ \\
Ceros de $\Xm/\Theta$ & $s = \pm j\,\SI{22.37}{\radian\per\second}$ \\
Pendiente servo $\rightarrow$ plato (diseño) & \SI{0.343}{\degree\per\degree} \\
\bottomrule
\end{tabular}
\end{table}

El doble polo en el origen tiene tres consecuencias que salen del modelo. Primero, la planta no
es estable en lazo abierto: con \SI{2}{\degree} de inclinación la bola llega a la pared en
aproximadamente \SI{1.0}{\second}, de modo que un ensayo escalón en lazo abierto no es posible.
Segundo, el método SIMC usado en el Proyecto~1 no aplica, porque supone una planta estable de
primer orden con tiempo muerto \cite{skogestad2003}. Tercero, un PI no puede estabilizar un doble
integrador: hace falta acción derivativa o realimentación de velocidad.

Para cuantificar el error de la aproximación lineal no se pueden usar escalones, porque la planta
no tiene equilibrio estable. Se aplicó un triplete de inclinación ($+A$, $-A$, $+A$) que, en el
modelo lineal, deja la bola en reposo en el punto de partida, con transiciones rápidas (el plato
a la velocidad máxima del servo, unos \SI{206}{\degree\per\second}) y lentas. La
Fig.~\ref{fig:barrido} muestra el error de trayectoria en función de la amplitud. Con el criterio
fijado por el equipo, un error menor al \SI{5}{\percent} de $\max|x|$, el modelo lineal vale
hasta $A = \SI{3}{\degree}$ con transiciones lentas, pero solo hasta \SI{0.5}{\degree} con
transiciones rápidas. El término dominante no es el seno, que a \SI{13}{\degree} equivale a solo
\SI{0.11}{\degree} de error de inclinación, sino el centrífugo $x\,\dot\theta^{2}$, que crece con
el cuadrado de la velocidad del plato. Por eso conviene que el controlador no pida cambios
bruscos de inclinación.

\begin{figure}[!t]
\centering
\begin{tikzpicture}
\begin{axis}[estilo, height=3.8cm, xlabel={Amplitud del triplete $A$ (\si{\degree})},
             ylabel={Error de trayectoria (\si{\percent})}, ymin=0, ymax=100, xmin=0, xmax=13.5,
             legend pos=north west]
  \addplot[cRojo, mark=*, mark size=1pt]
    table[col sep=comma, x=A_deg, y=err_rel_x_pct,
          restrict expr to domain={\thisrow{modo_1rapida_2lenta}}{1:1}] {\G/A2_barrido_amplitud.csv};
  \addplot[cAzul, mark=square*, mark size=1pt]
    table[col sep=comma, x=A_deg, y=err_rel_x_pct,
          restrict expr to domain={\thisrow{modo_1rapida_2lenta}}{2:2}] {\G/A2_barrido_amplitud.csv};
  \addplot[cGris, dashed, domain=0:13.5, samples=2] {5};
  \legend{Transición rápida, Transición lenta, Umbral \SI{5}{\percent}}
\end{axis}
\end{tikzpicture}
\caption{Error del modelo lineal frente al no lineal, en porcentaje de $\max|x|$, para el
triplete de inclinación con transiciones rápidas y lentas.}
\label{fig:barrido}
\end{figure}

\subsection{Caracterización del sensor y del actuador}

Antes de identificar la dinámica de la bola se midieron, sobre el prototipo, los elementos de la
cadena que el modelo trataba como ideales. La Tabla~\ref{tab:caracterizacion} resume las
mediciones.

\begin{table}[!t]
\caption{Mediciones de la cadena de medición y del actuador sobre el prototipo. Fuente:
\texttt{datos\_prototipo/camara/} y \texttt{datos\_prototipo/mecanica/}.}
\label{tab:caracterizacion}
\centering
\footnotesize
\begin{tabular}{@{}lr@{}}
\toprule
\textbf{Magnitud} & \textbf{Valor medido} \\
\midrule
Periodo entre cuadros (YUYV, 2 búferes) & \SI{33.25}{\milli\second}, desv. \SI{1.9}{\milli\second} \\
Cuadros dentro de $\pm\SI{5}{\milli\second}$ / perdidos & \SI{100}{\percent} / \SI{0}{\percent} \\
Latencia de la visión (mediana, 2 corridas) & \SI{57.8}{} y \SI{60.3}{\milli\second} \\
Ruido de posición, bola quieta (centro / esquina) & \SI{0.12}{} / \SI{0.09}{\milli\metre} \\
Pendiente servo $\rightarrow$ plato cerca del neutro (v3) & \SI{0.47}{\degree\per\degree} \\
Holgura del mecanismo (eje X / eje Y) & hasta \SI{0.8}{} / \SI{2.4}{\degree} \\
\bottomrule
\end{tabular}
\end{table}

El periodo de muestreo quedó en $\Ts = \SI{33.3}{\milli\second}$ con el formato YUYV y dos
búferes; el MJPG se descartó porque perdía el \SI{46}{\percent} de los cuadros. La latencia
medida, de unos \SI{60}{\milli\second}, quedó por debajo de los \SI{75}{\milli\second} que se
habían estimado para el diseño. La relación entre el ángulo del servo y el del plato resultó
asimétrica y con una pendiente cerca de 1,4 veces la del diseño, así que el firmware convierte
la inclinación pedida en comando del servo con la tabla medida y no con la cinemática del diseño.
Con esa tabla, el peor eje alcanza unos $\pm\SI{11.3}{\degree}$, y la saturación del controlador
se fijó en $\thlim = \SI{10}{\degree}$. La tabla vigente es la v3, medida el 2026-11-09 después
de reajustar los brazos de los servos.

\subsection{Identificación experimental}

Como la planta no es estable en lazo abierto, la identificación se hizo en lazo cerrado. El
firmware de prueba de la ESP32 ejecuta un ensayo automático con escalones de referencia, una
perturbación pseudoaleatoria de $\pm\SI{1.5}{\degree}$ sumada a la inclinación y trayectorias
circulares, y devuelve en cada cuadro la inclinación que aplicó. Con esa telemetría y la posición
medida se ajustó, por mínimos cuadrados y en ventanas de \SI{1}{\second}, el modelo por eje

\begin{equation}
\ddot x = -c_{\mathrm{ef}}\,g\,\sin\theta(t-\tau) + b - f\,\operatorname{sign}(\dot x),
\label{eq:b3}
\end{equation}

donde $c_{\mathrm{ef}}$ es una ganancia de rodadura efectiva, que incluye también la ganancia
real de la transmisión servo--plato, $b$ representa un desnivel constante, $f$ un rozamiento o
zona muerta, y $\tau$ el retardo entre la orden de inclinación y la respuesta de la bola. El
retardo se buscó en una grilla, minimizando el error de ajuste. En la Fig.~\ref{fig:tau}, el
mínimo está en \SI{340}{\milli\second} para X y en \SI{270}{\milli\second} para Y, y la suma de
los dos ejes es mínima en \SI{270}{\milli\second}, que se tomó como retardo común. Las curvas son
bastante planas alrededor del mínimo, así que el retardo queda determinado con poca precisión.

\begin{figure}[!t]
\centering
\begin{tikzpicture}
\begin{axis}[estilo, height=3.6cm, xlabel={Retardo supuesto $\tau$ (\si{\second})},
             ylabel={SSE (\si{\milli\metre\squared})}, legend pos=north east,
             scaled y ticks=base 10:-3, ytick scale label code/.code={}, ylabel={SSE ($\times 10^{3}$ \si{\milli\metre\squared})}]
  \addplot[cAzul]  table[col sep=comma, x=tau_s, y=sse_x] {\G/B3_tau.csv};
  \addplot[cNaranja, dashed] table[col sep=comma, x=tau_s, y=sse_y] {\G/B3_tau.csv};
  \legend{Eje X, Eje Y}
\end{axis}
\end{tikzpicture}
\caption{Error de ajuste del modelo \eqref{eq:b3} en función del retardo supuesto, ensayo
\texttt{sesion\_20261113\_151156}.}
\label{fig:tau}
\end{figure}

La Tabla~\ref{tab:b3} reúne las identificaciones hechas a medida que cambiaba el prototipo.
Salvo en el ensayo inmediatamente posterior a la lubricación, $c_{\mathrm{ef}}$ salió entre
\num{0.12} y \num{0.28}, de tres a seis veces menos que el
\num{0.714} del modelo, y el retardo entre \SI{0.27}{} y \SI{0.41}{\second}, cuatro a cinco
veces el estimado. Al cambiar la bola de PLA por una pelota de ping pong, que es lisa y liviana,
se obtuvieron el mismo $c_{\mathrm{ef}}$ y la misma zona muerta de alrededor de \SI{1}{\degree}.
Eso indica que el efecto no viene de la bola sino del mecanismo: holgura de rótulas y varillas, y
respuesta del servo a cambios pequeños, que quedan dentro de la perturbación de
$\pm\SI{1.5}{\degree}$ del ensayo.

\begin{table}[!t]
\caption{Parámetros identificados en lazo cerrado (B3) por ensayo, en el formato eje X / eje
Y. $\tau$ es el óptimo de cada eje. ZM: zona muerta equivalente, en grados pedidos. ---: no
registrado. Fuente:
\texttt{estado\_proyecto.md} y \texttt{etapaB/resultados/}.}
\label{tab:b3}
\centering
\footnotesize
\setlength{\tabcolsep}{3pt}
\begin{tabular}{@{}llccc@{}}
\toprule
\textbf{Ensayo} & \textbf{Condición} & $c_{\mathrm{ef}}$ & $\tau$ (\si{\milli\second}) & ZM (\si{\degree}) \\
\midrule
\texttt{143311} & PLA, tabla v2      & 0,19 / 0,22 & 340 / 370 & 1,2 / 1,0 \\
\texttt{164915} & PLA, tabla v3      & 0,16 / 0,24 & 340 / 310 & 1,0 / 1,0 \\
\texttt{165209} & PLA, tabla v3      & 0,12 / 0,21 & 270 / 370 & 1,9 / 1,5 \\
\texttt{105834} & Ping pong          & 0,19 / 0,22 & 270 / 410 & 1,2 / 1,0 \\
\texttt{150437} & Tras lubricar      & 0,06 / 0,24 & --- & 2,5 / --- \\
\texttt{151156} & Tras lubricar      & 0,18 / 0,28 & 340 / 270 & 0,86 / 1,18 \\
\bottomrule
\end{tabular}
\end{table}

El ensayo \texttt{150437} muestra además que el mecanismo no es estable en el tiempo: después
de lubricar el servo X, la ganancia efectiva de ese eje cayó a \num{0.06}, y en el ensayo
siguiente se recuperó, pero con el neutro corrido unos \SI{-1.7}{\degree}.

\subsection{Validación y selección del modelo}

Cada modelo identificado se validó con un ensayo distinto del usado para ajustarlo: desde la
posición y la velocidad de los primeros seis cuadros de cada ventana se predijeron los
\SI{0.8}{\second} siguientes sin reajustar parámetros. La Tabla~\ref{tab:validacion} compara el
modelo identificado con el nominal de la Etapa~A para el ensayo más reciente.

\begin{table}[!t]
\caption{Validación del modelo identificado con \texttt{151156} sobre los datos de
\texttt{112242}, frente al modelo nominal ($c = 5/7$, $b = 0$, $f = 0$,
$\tau = \SI{75}{\milli\second}$). Predicción de
\SI{0.8}{\second} por ventana. Fuente: \texttt{etapaB/resultados/B3\_consola.txt}.}
\label{tab:validacion}
\centering
\footnotesize
\setlength{\tabcolsep}{3pt}
\begin{tabular}{@{}llS[table-format=2.1]S[table-format=2.1]S[table-format=+2.1]S[table-format=1.3]S[table-format=2.1]S[table-format=1.2]@{}}
\toprule
\textbf{Modelo} & \textbf{Eje} & {RMSE} & {MAE} & {Sesgo} & {$R^2$} & {FIT} & {$\rho_1$} \\
 & & {(\si{\milli\metre})} & {(\si{\milli\metre})} & {(\si{\milli\metre})} & & {(\si{\percent})} & \\
\midrule
Identificado & X & 15.0 & 9.9  & +8.3  & 0.906 & 69.3 & 0.89 \\
Identificado & Y & 13.7 & 9.1  & +7.0  & 0.904 & 69.1 & 0.89 \\
Nominal      & X & 16.5 & 9.8  & -2.1  & 0.887 & 66.4 & 0.91 \\
Nominal      & Y & 21.9 & 14.9 & -12.1 & 0.755 & 50.5 & 0.88 \\
\bottomrule
\end{tabular}
\end{table}

El modelo identificado predice mejor en los dos ejes, sobre todo en Y, donde el FIT sube de
\SI{50.5}{} a \SI{69.1}{\percent}. Sin embargo, el residuo tiene una autocorrelación a un paso
$\rho_1$ de \num{0.89}: no es ruido blanco, sino un error con estructura, de modo que al modelo
\eqref{eq:b3} todavía le falta una parte de la dinámica. Lo más probable es que sea la holgura
del mecanismo y la dinámica real del servo, que el modelo reduce a una ganancia, un retardo y
una zona muerta. Además, el $c_{\mathrm{ef}}$ de unos \num{0.2} se obtiene con perturbaciones
pequeñas y rápidas; el balance de fuerzas en las trayectorias circulares, que son movimientos
lentos y grandes, dio alrededor de \num{0.57}. Esa diferencia es coherente con una transmisión que
pierde los movimientos pequeños dentro del juego mecánico.

Se seleccionó como modelo nominal el identificado con el ensayo \texttt{151156}, por tres
razones. Corresponde al estado del prototipo con el que se hicieron las últimas pruebas, con la
pelota de ping pong y los servos después de la lubricación. Fue validado con un ensayo
independiente, donde predijo mejor que el modelo de la Etapa~A en los dos ejes. Y sus parámetros
salen de una sola identificación, sin promediar modelos de ensayos distintos. Con el retardo
común de \SI{270}{\milli\second} que eligió la grilla, el modelo nominal por eje es

\begin{equation}
G_{p,\mathrm{nom}}(s) = \frac{X(s)}{\Theta(s)} = -\frac{c_{\mathrm{ef}}\,g\;e^{-\tau_{\mathrm{nom}}s}}{s^{2}},
\label{eq:gpnom}
\end{equation}

con $c_{\mathrm{ef}}\,g = \SI{1780}{\milli\metre\per\second\squared\per\radian}$ en X
($c_{\mathrm{ef}} = \num{0.18}$), \SI{2765}{\milli\metre\per\second\squared\per\radian} en Y
($c_{\mathrm{ef}} = \num{0.28}$) y $\tau_{\mathrm{nom}} = \SI{0.27}{\second}$, más una zona
muerta de \SI{0.86}{\degree} en X y \SI{1.18}{\degree} en Y. El desnivel identificado se corrige
en el firmware con un ajuste fijo de la inclinación (\emph{trim}) y no forma parte del modelo.
La Fig.~\ref{fig:validacion} muestra algunas ventanas de la validación. El modelo nominal sigue
la tendencia de la bola, pero en varias ventanas se separa de la medición de forma sostenida,
que es lo que indica la autocorrelación del residuo.

\begin{figure}[!t]
\centering
\begin{tikzpicture}
\begin{groupplot}[group style={group size=1 by 2, vertical sep=0.6cm}, estilo, height=3.3cm,
                  xmin=0, xmax=155, xticklabels={}, unbounded coords=jump]
\nextgroupplot[ylabel={$x$ (\si{\milli\metre})}, legend pos=north east, legend columns=3]
  \addplot[cAzul, line width=0.9pt, restrict expr to domain={\thisrow{eje}}{1:1}]
    table[col sep=comma, x=k, y=x_med_mm] {\G/B3_validacion_ventanas.csv};
  \addplot[cVerde, dashed, restrict expr to domain={\thisrow{eje}}{1:1}]
    table[col sep=comma, x=k, y=x_ident_mm] {\G/B3_validacion_ventanas.csv};
  \addplot[cRojo, densely dotted, restrict expr to domain={\thisrow{eje}}{1:1}]
    table[col sep=comma, x=k, y=x_nom_mm] {\G/B3_validacion_ventanas.csv};
  \legend{Medida, Identificado, Etapa~A}
\nextgroupplot[ylabel={Residuo (\si{\milli\metre})}, xlabel={Muestra (6 ventanas de \SI{0.8}{\second})}]
  \addplot[cVerde, restrict expr to domain={\thisrow{eje}}{1:1}]
    table[col sep=comma, x=k, y=res_ident_mm] {\G/B3_validacion_ventanas.csv};
  \addplot[cGris, dashed, domain=0:155, samples=2] {0};
\end{groupplot}
\end{tikzpicture}
\caption{Validación del modelo nominal en el eje X: primeras seis ventanas del ensayo
\texttt{112242}, que no se usó para identificar. Arriba, medición y predicciones; abajo, residuo
del modelo identificado.}
\label{fig:validacion}
\end{figure}
